\def\probtitle{크림 먼저, 잼 먼저}
\def\probno{C} % 문제 번호

\begin{problem}{\probtitle}{표준 입력(stdin)}{표준 출력(stdout)}{1 초}{1024 megabytes}

영국에서 스콘을 맛있게 먹는 방법은 위에 크림을 바른 후 잼을 얹는 것, 혹은 위에 잼을 바른 후 크림을 얹는 것이다. 성준이는 파티에서 어느 방식의 스콘이 더 많은지 세어 보기로 했다. 불행하게도, 모든 스콘들은 하나의 탑으로 쌓여 있어, 위에서부터 재료를 읽을 수밖에 없었다.

스콘 탑의 재료들을 나타내는 문자열 $T$가 주어지면, \textbf{크림이 위인 스콘의 개수}와, \textbf{잼이 위인 스콘의 개수}를 각각 구하여라.

문자열 $T$는 문자 \texttt{C}, \texttt{J}, \texttt{S} 로만 이루어져 있다. \texttt{C}는 크림을, \texttt{J}는 잼을, \texttt{S}는 스콘을 나타낸다. 

\begin{itemize}[topsep=0pt,noitemsep]
  \item \textbf{크림이 위인 스콘}: 1개 이상의 \texttt{C}, 1개 이상의 \texttt{J}, 하나의 \texttt{S}가 차례대로 연속해서 나타나는 경우이다.
  \item \textbf{잼이 위인 스콘}: 1개 이상의 \texttt{J}, 1개 이상의 \texttt{C}, 하나의 \texttt{S}가 차례대로 연속해서 나타나는 경우이다.
  \item \texttt{C}와 \texttt{J}는 모두 자신의 오른쪽에서 가장 처음 오는 \texttt{S}에 바르는 것으로 가정한다. 즉, \texttt{CJCS}에서 \texttt{S}에는 \texttt{C}, \texttt{J}, \texttt{C}가 모두 발라져 있다.
\end{itemize}

만약 크림이나 잼 중 하나라도 빠져 있거나, 두 재료가 번갈아 여러 번 나타나는 스콘은 잘못된 스콘으로 간주하며, \textbf{어느 쪽에도 세지 않는다.}

\InputFile

첫째 줄에 문자열 $T$의 길이를 나타내는 정수 $N$이 주어진다.

둘째 줄에 길이가 $N$인 문자열 $T$가 주어진다.

\OutputFile

첫째 줄에 크림이 위인 스콘의 개수를 출력한다.

둘째 줄에 잼이 위인 스콘의 개수를 출력한다.

\Constraints

\begin{itemize}[topsep=0pt,noitemsep]
    \item $1\leq N\leq 200$
    \item 문자열 $T$는 \texttt{C, J, S} 로만 구성되어 있으며, 반드시 \texttt{S}로 끝난다.
\end{itemize}

\Example

\begin{example}
    \exmpfile{./example/01.in.txt}{./example/01.out.txt}%
    \exmpfile{./example/02.in.txt}{./example/02.out.txt}%
    \exmpfile{./example/03.in.txt}{./example/03.out.txt}%
    \exmpfile{./example/04.in.txt}{./example/04.out.txt}%
\end{example}

\Notes

예제 4의 문자열 $T=\texttt{SJJJCSCJSJSCJCSJJCCJJJCJSSJS}$ 를 \texttt{S}를 기준으로 나누면
\[
\texttt{S},\ \texttt{JJJCS},\ \texttt{CJS},\ \texttt{JS},\ \texttt{CJCS},\ \texttt{JJCCJJJCJS},\ \texttt{S},\ \texttt{JS}
\]
를 얻는다.

이 중 \texttt{CJS}는 크림이 위인 스콘 1개, \texttt{JJJCS}는 잼이 위인 스콘 1개이다.
나머지 \texttt{S}, \texttt{JS}, \texttt{CJCS}, \texttt{JJCCJJJCJS}, \texttt{S}, \texttt{JS}는
재료가 빠져 있거나 재료가 여러 번 번갈아 나타나므로 모두 잘못된 스콘이다.

따라서 출력은 첫째 줄에 $1$, 둘째 줄에 $1$이다.

\end{problem}